\section{总结与延伸}

本节把 EP116 压缩成三个层次。第一，公司史层面：明略从推荐系统、秒针、明略数据、并购和痛苦急转中，积累了 ToB 数据、交付和经营能力。第二，技术商业层面：企业级 AI 的差异化来自私有数据、工作流、Tool Use、Agentic Model 和可信交付，而不是公开 token 的价格战。第三，组织层面：AI 会重写供给侧能力、链接网络、岗位结构和 founder 的使命表达。

\begin{importantbox}{把 EP116 放进张小珺 AI 队列}
EP117 讲论文和技术史，EP118 讲 VLA、Agent OS 和平台，EP116 则讲企业服务场景里的 AI 落地：它提醒我们，AI 的产业价值不只在模型公司，也在拥有私有数据、行业流程和交付责任的 ToB 公司。
\end{importantbox}

\subsection{关键 takeaways}

\begin{enumerate}
\item 企业级 AI 的护城河不是公开模型，而是私有数据、流程和可信交付。
\item Agentic Model 在 ToB 里必须接入工具、权限、审计和真实 workflow。
\item 只提供链接网络、不掌握供给侧能力的公司，会被 AI 重新定价。
\item ToB 创业的核心难点是现金流、交付、聚焦和组织耐力，不是单点技术。
\item Trusted Agentic Model 比“最聪明模型”更贴近企业客户的真实购买理由。
\end{enumerate}

\subsection{拓展阅读}

\begin{itemize}
\item 对 Agentic Model 感兴趣，可对照 EP139 Agent 综述和 EP138 罗福莉访谈，理解 Agent 从技术 loop 到组织/后训练体系的迁移。
\item 对企业私有数据感兴趣，可对照 EP134 数据综述，理解数据定价、Recipe 和真实部署数据的差异。
\item 对 AI 平台和操作系统感兴趣，可对照 EP118 李想访谈与 EP127 AI War，观察 Agent OS、企业 workflow 和平台入口之间的关系。
\end{itemize}

\end{document}
